\documentclass[11pt]{article}
\usepackage[T1]{fontenc}
\usepackage{textcomp}
\usepackage[margin=0.75in]{geometry}
\usepackage{amsmath,amssymb,amsthm}
\usepackage{array,booktabs}
\usepackage{hyperref}
\setlength{\parindent}{0pt}
\newtheorem{theorem}{Theorem}
\theoremstyle{definition}
\newtheorem{definition}{Definition}
\title{Flow Proof Artifact: Real-add-cancel}
\author{Generated by \texttt{flow doc proof}}
\date{Flow Proof Book}
\begin{document}
\maketitle
Equal reals can be subtracted from both sides.\\[0.5em]
\textbf{Source.} landau — *Foundations of Analysis*\\[0.5em]
\bigskip
\noindent{\large \textbf{Derived fact 1.} \textit{x + z = y + z implies x = y} \hfill the real numbers \cdot addition \cdot right cancellation holds}\par
\smallskip\noindent\textit{Coordinate: the real numbers \cdot addition \cdot right cancellation holds}\par
\smallskip\noindent\textit{Source: landau}\par
\smallskip\noindent\textit{Built on: parentheses do not matter, for addition on the real numbers, zero is the right identity, for addition on the real numbers}\par
\medskip
\noindent\textbf{Goal.} x + z = y + z implies x = y\par
\noindent$\displaystyle \forall x \in \mathbb{R} \forall y \in \mathbb{R} \forall z \in \mathbb{R}\quad x = y$\par
\medskip
\noindent
\begin{tabular}{@{} >{\raggedright\arraybackslash}p{0.06\textwidth} >{\raggedright\arraybackslash}p{0.47\textwidth} >{\raggedleft\arraybackslash}p{0.41\textwidth} @{}}
\textbf{\#} & \textbf{Proof} & \textbf{Mathematics} \\
\midrule
\textbf{1.} & We prove that right cancellation holds for addition on the real numbers. & \\[0.45em]
\textbf{2.} & We split into exhaustive cases — the claim must hold in each one. & \\[0.45em]
\textbf{3.} & Case 1 (see step 2): suppose x + z  equals  y + z. & \\[0.45em]
\textbf{4.} & We invoke the derived fact governing addition on the real numbers: parentheses do not matter, for addition on the real numbers (instantiated for x, z, -z). & \\[0.45em]
\textbf{5.} & We invoke the definitional clause governing addition on the real numbers: zero is the right identity, for addition on the real numbers (instantiated for x). & \\[0.45em]
\textbf{6.} & From step 3, step 4, and step 5, this implies x equals y. Together with the other cases (step 3), the goal is discharged. Hence proven. & $\displaystyle x = y$ \\[0.45em]
\bottomrule
\end{tabular}
\label{thm:Real-addition-right_cancellation_holds}
\par\medskip\hrule
\end{document}
